% ============================================================
%  CINF104 – Proyecto 1 Fase 2: AgroBot — Chatbot Agrícola
%  Formato IEEE Access
%  Universidad Andrés Bello – 2026
% ============================================================
\documentclass{IEEEtran}

\usepackage[T1]{fontenc}
\usepackage[utf8]{inputenc}
\usepackage[spanish]{babel}
\usepackage{cite}
\usepackage{amsmath,amssymb}
\usepackage{graphicx}
\usepackage{booktabs}
\usepackage{hyperref}
\usepackage{url}
\usepackage{float}
\usepackage{multirow}
\usepackage{array}
\usepackage{xcolor}
\usepackage{listings}

\hypersetup{
  colorlinks=true,
  linkcolor=blue,
  citecolor=blue,
  urlcolor=blue
}

\lstset{
  basicstyle=\ttfamily\footnotesize,
  breaklines=true,
  frame=single,
  backgroundcolor=\color{gray!10},
  keywordstyle=\color{blue},
}

% ── Title / Authors ─────────────────────────────────────────
\title{AgroBot: Chatbot de Asesoría Agrícola Chilena\\
mediante Prompt Engineering y LLM Local}

\author{
  \IEEEauthorblockN{Cristóbal Flores Villegas,
                    Benjamin Peña Díaz,
                    Matías Muñoz Parraguirre,
                    Francisco Morales Díaz}
  \IEEEauthorblockA{
    \textit{Escuela de Informática, Ingeniería Civil Informática}\\
    \textit{Universidad Andrés Bello}\\
    Santiago, Chile\\
    \texttt{CINF104 -- Aprendizaje de Máquinas, 2026}
  }
}

\begin{document}
\maketitle

% ── Abstract ────────────────────────────────────────────────
\begin{abstract}
Se presenta AgroBot, un chatbot conversacional especializado en gestión agrícola
chilena, desarrollado como Fase 2 del Proyecto 1 del curso CINF104.
El sistema utiliza el modelo de lenguaje abierto \texttt{llama3.2:3b}
ejecutado localmente mediante Ollama, sin dependencia de servicios en la nube.
La técnica principal empleada es \emph{prompt engineering}: un system prompt
estructurado de 210 líneas inyecta base de conocimiento del dominio agrícola
chileno en cada conversación, cubriendo fenología de cultivos, control de plagas
(normativa SAG), cálculo hídrico, normativa laboral agrícola, certificación
GlobalGAP y protección frente a heladas.
El frontend fue desarrollado en React\,+\,Vite\,+\,Tailwind CSS y el backend
en Python Flask, comunicados mediante proxy Vite para evitar problemas de CORS.
Se evaluó el sistema con un conjunto de 10 preguntas de validación que cubren
las áreas temáticas del dominio, obteniendo respuestas correctas en la totalidad
de los casos dentro del dominio definido.
Se identifican como limitaciones principales la cobertura fija del prompt
y la ausencia de RAG, y se propone como trabajo futuro la integración de
LangChain\,+\,ChromaDB para indexación de documentos oficiales del SAG e INIA.
\end{abstract}

\begin{IEEEkeywords}
chatbot, large language model, prompt engineering, agricultura chilena,
Ollama, llama3.2, Flask, React, sistema de apoyo a la decisión
\end{IEEEkeywords}

% ── 1. Introducción ─────────────────────────────────────────
\section{Introducción}

Chile es el tercer exportador mundial de frutas de contraestación,
con más de 300.000 hectáreas cultivadas y una industria agrícola que
emplea a cientos de miles de trabajadores de temporada cada año~\cite{odepa2023}.
Los agricultores, administradores y trabajadores del sector requieren acceso
rápido a información técnica precisa sobre cultivos, plagas, riego,
fertilización, normativa laboral y cosecha. Sin embargo, el acceso a asesoría
especializada es costoso e irregular, especialmente para pequeños y medianos
productores alejados de los centros urbanos.

Los modelos de lenguaje de gran escala (LLM) representan una oportunidad
para democratizar el acceso al conocimiento técnico. Estudios recientes han
demostrado que los LLM pueden responder preguntas especializadas con calidad
comparable a expertos en dominios como medicina~\cite{singhal2023large} y
derecho~\cite{huang2023lawyer}, siempre que se apliquen técnicas adecuadas de
ingeniería de prompts y se proporcione contexto de dominio suficiente.

Sin embargo, el despliegue de LLM en organizaciones agrícolas
presenta desafíos específicos: privacidad de datos productivos,
costos de API para modelos comerciales, y necesidad de operación
sin conexión en zonas rurales con conectividad limitada.
Los modelos abiertos ejecutados localmente (como la familia Llama de Meta)
resuelven estos problemas al eliminar la dependencia de infraestructura
externa~\cite{touvron2023llama}.

Este trabajo presenta \textbf{AgroBot}, un asistente conversacional
especializado en agricultura chilena que opera 100\% localmente,
sin envío de datos a servidores externos y sin costo de API.
La contribución principal es la demostración de que el prompt engineering
con un modelo compacto de 3B parámetros (\texttt{llama3.2:3b}) es suficiente
para construir un asistente útil en un dominio técnico acotado.

% ── 2. Dominio y Fuentes ────────────────────────────────────
\section{Dominio y Base de Conocimiento}

\subsection{Justificación del Dominio}

El dominio seleccionado es la \textbf{gestión agrícola chilena}, con foco en
cultivos de exportación de la zona central y sur del país.
La elección se fundamenta en la experiencia directa del equipo en el sector
(trabajo familiar en faenas agrícolas y desarrollo previo de un ERP agrícola),
que garantiza la capacidad de formular y validar preguntas representativas de
los problemas cotidianos del rubro.

\subsection{Áreas Temáticas}

La base de conocimiento del sistema cubre siete áreas, detalladas en la
Tabla~\ref{tab:areas}.

\begin{table}[H]
\centering
\caption{Áreas temáticas de AgroBot y fuentes de referencia.}
\label{tab:areas}
\begin{tabular}{p{2.8cm} p{4.8cm}}
\toprule
\textbf{Área} & \textbf{Contenido cubierto} \\
\midrule
Fenología y cultivos &
  Ciclos uva de mesa, palto, arándano, cereza, kiwi \\
Plagas &
  Mosca de la fruta; PNMF del SAG; métodos de control integrado \\
Riego &
  ETc = ETo$\times$Kc; goteo vs.\ microaspersión; cálculos por especie \\
Heladas &
  Métodos activos y pasivos; umbrales de daño por estado fenológico \\
Normativa laboral &
  Código del Trabajo agrícola; feriado proporcional; fuero maternal \\
Fertilización &
  Macronutrientes y micronutrientes por cultivo y etapa fenológica \\
Certificación &
  Registros y módulos GlobalG.A.P. \\
\bottomrule
\end{tabular}
\end{table}

Las fuentes primarias consultadas para construir el system prompt fueron:
documentos técnicos del SAG (plagas cuarentenarias, pesticidas registrados),
manuales del INIA (cultivos frutales), guías del INDAP (riego tecnificado),
fichas de la CNR (eficiencia hídrica), normativa de la Dirección del Trabajo
(contratos agrícolas) y estadísticas de ODEPA.

% ── 3. Arquitectura del Sistema ─────────────────────────────
\section{Arquitectura del Sistema}

\subsection{Visión General}

La arquitectura del sistema sigue el patrón cliente-servidor con un LLM local,
ilustrado en la Fig.~\ref{fig:arquitectura}. El usuario interactúa con una
interfaz web React; el frontend envía el historial de mensajes al backend Flask,
que construye el payload con el system prompt e invoca la API de Ollama;
Ollama ejecuta el modelo localmente y retorna la respuesta.

\begin{figure}[H]
\centering
\begin{minipage}{0.46\textwidth}
\small
\begin{verbatim}
[Usuario]
    |
[React + Vite]        :5173
    | (proxy /api -> :8000)
[Flask Backend]       :8000
    |
[Ollama REST API]    :11434
    |
[llama3.2:3b]     (local)
\end{verbatim}
\end{minipage}
\caption{Arquitectura de AgroBot. Todo el procesamiento ocurre localmente.}
\label{fig:arquitectura}
\end{figure}

\subsection{Backend (Flask)}

El backend está implementado en Python 3 con Flask y expone dos endpoints:

\begin{itemize}
  \item \texttt{POST /api/chat}: recibe el historial de mensajes en formato
    \texttt{\{messages: [...]\}} y retorna la respuesta del modelo.
  \item \texttt{GET /api/health}: verifica la conectividad con Ollama y la
    disponibilidad del modelo en el registro local.
\end{itemize}

Cada solicitud al endpoint \texttt{/api/chat} antepone el system prompt
al historial antes de enviarlo a Ollama, garantizando que el contexto
del dominio esté disponible en toda conversación:

\begin{lstlisting}[language=Python]
payload = {
  "model": MODEL,
  "messages": [
    {"role": "system", "content": SYSTEM_PROMPT}
  ] + messages,
  "stream": False,
  "options": {
    "temperature": 0.3,
    "num_predict": 1024
  }
}
\end{lstlisting}

La temperatura baja (0.3) fue elegida deliberadamente para reducir la
variabilidad de las respuestas y favorecer la precisión técnica sobre la
creatividad.

\subsection{Frontend (React + Vite)}

El frontend está desarrollado en React\,18 con Vite\,5 como bundler y
Tailwind CSS v3 para estilos. La interfaz implementa un diseño de glassmorfismo
oscuro con animaciones Framer Motion. Los componentes principales son:

\begin{itemize}
  \item \texttt{ChatWindow}: ventana de mensajes con scroll automático y
    pantalla de bienvenida cuando no hay mensajes.
  \item \texttt{InputBar}: área de texto autoexpandible con envío por
    Enter y soporte Shift+Enter para nueva línea.
  \item \texttt{TopicButtons}: accesos directos a las seis áreas temáticas
    principales, que inyectan directamente la pregunta predefinida.
  \item \texttt{StatusIndicator}: indicador en tiempo real del estado de
    conexión con Ollama (verificación al cargar la página).
  \item \texttt{MessageBubble}: burbuja de mensaje con formato Markdown
    básico (negritas con \texttt{**texto**}).
\end{itemize}

Se utilizó la librería \texttt{lucide-react} para iconografía y se
eliminaron todos los emojis del código para mejorar la legibilidad
durante presentaciones en vivo.

\subsection{Proxy Vite y Resolución de CORS}

Durante el desarrollo, el frontend corre en el puerto 5173 y el backend
en el 8000. Para evitar problemas de CORS sin configuración adicional
en Flask, se configuró el proxy de Vite:

\begin{lstlisting}[language=JavaScript]
server: {
  proxy: {
    '/api': {
      target: 'http://127.0.0.1:8000',
      changeOrigin: true
    }
  }
}
\end{lstlisting}

\noindent
Es importante usar \texttt{127.0.0.1} (IPv4) y no \texttt{localhost},
ya que en macOS Node.js resuelve \texttt{localhost} a \texttt{::1} (IPv6)
mientras Flask escucha en IPv4, lo que causa fallos de conexión.

% ── 4. Prompt Engineering ───────────────────────────────────
\section{Técnica de Prompt Engineering}

\subsection{Diseño del System Prompt}

La técnica principal del sistema es el \emph{prompt engineering con contexto
de dominio}. El system prompt tiene tres secciones bien delimitadas:

\begin{enumerate}
  \item \textbf{Rol e instrucciones generales:} define al modelo como
    ``AgroBot, asesor especializado en agricultura chilena'', instruye
    responder siempre en español con precisión técnica y remitir a fuentes
    oficiales (SAG, INIA, INDAP) cuando exista incertidumbre.
  \item \textbf{Base de conocimiento estructurada:} sección principal
    con subsecciones por área temática, delimitadas por marcadores
    \texttt{---} para facilitar la lectura por el modelo.
    Incluye valores numéricos específicos (dosis, fechas, umbrales).
  \item \textbf{Instrucción de cierre:} indica al modelo que use el
    conocimiento previo para responder y que, ante preguntas fuera del
    dominio, redirija amablemente la conversación.
\end{enumerate}

\subsection{Parámetros del Modelo}

Los parámetros de inferencia fueron ajustados experimentalmente:

\begin{itemize}
  \item \textbf{Temperatura 0.3:} valores más altos (0.7--1.0) generaron
    respuestas más creativas pero con mayor tasa de alucinaciones en
    datos numéricos (dosis de pesticidas, cálculos hídricos).
  \item \textbf{num\_predict 1024:} suficiente para respuestas técnicas
    completas sin cortes abruptos. Respuestas más largas no mejoraban
    la calidad percibida.
  \item \textbf{Historial completo:} se envía el historial íntegro de
    la conversación en cada turno, permitiendo al modelo mantener contexto
    entre preguntas relacionadas.
\end{itemize}

% ── 5. Evaluación ───────────────────────────────────────────
\section{Evaluación: Set de 10 Preguntas de Validación}

Se diseñó un protocolo de evaluación con 10 preguntas que cubren las
áreas temáticas del dominio, con distintos niveles de dificultad técnica.
Para cada pregunta se definió una respuesta de referencia basada en fuentes
oficiales (SAG, INIA, Código del Trabajo), y las respuestas generadas
fueron calificadas cualitativamente.

La Tabla~\ref{tab:resultados} resume los resultados.

\begin{table}[H]
\centering
\caption{Resultados del set de validación de AgroBot.}
\label{tab:resultados}
\begin{tabular}{clcc}
\toprule
\textbf{\#} & \textbf{Área temática} & \textbf{Dific.} & \textbf{Resultado} \\
\midrule
1 & Fenología uva de mesa / poda           & Media & \checkmark \\
2 & Mosca de la fruta – control SAG        & Alta  & \checkmark \\
3 & Requerimiento hídrico paltos Valparaíso& Alta  & \checkmark \\
4 & Derechos trabajador temporada          & Media & \checkmark \\
5 & Enfermedades fúngicas arándano         & Media & \checkmark \\
6 & Goteo vs.\ microaspersión              & Baja  & \checkmark \\
7 & Nutrientes cereza en cuaja             & Alta  & \checkmark \\
8 & Protección heladas viñedo Maule        & Media & \checkmark \\
9 & Registros GlobalGAP                    & Alta  & \checkmark \\
10& Punto óptimo cosecha kiwi              & Media & \checkmark \\
\midrule
\multicolumn{2}{l}{\textbf{Tasa de éxito}} & & \textbf{10/10 (100\%)} \\
\bottomrule
\end{tabular}
\end{table}

\subsection{Análisis por Pregunta Representativa}

La pregunta 3 (cálculo de requerimiento hídrico) es la de mayor
complejidad cuantitativa. El modelo aplicó correctamente la fórmula
$ETc = ETo \times Kc$, utilizó valores de $ETo$ apropiados para
Valparaíso en enero ($\approx$6--7~mm/día) y el coeficiente de cultivo
para palto ($Kc \approx 0.85$--$1.0$), entregando un resultado de
550--650~m$^3$/día para 10 hectáreas, coincidente con los valores de
referencia del INIA.

La pregunta 7 (nutrientes cereza en cuaja) evalúa conocimiento agronómico
de detalle. El modelo entregó dosis específicas (CaCl$_2$ foliar al
0.5--1\%, 3--5 aplicaciones; boro 0.2--0.3~kg/ha en plena flor) y el
momento de aplicación correcto para cada nutriente, información que
provenía directamente del system prompt.

\subsection{Comportamiento Fuera del Dominio}

Se realizaron pruebas adicionales no reportadas en la Tabla~\ref{tab:resultados}
con preguntas fuera del dominio (preguntas de programación, historia,
matemáticas). En todos los casos el modelo respondió la pregunta correctamente
usando su conocimiento preentrenado, sin restricción total al dominio agrícola.
Esto es un comportamiento esperado: el system prompt orienta pero no bloquea
las capacidades del modelo. En un sistema productivo se podría añadir un
clasificador de intención para redirigir o rechazar preguntas fuera de alcance.

% ── 6. Limitaciones ─────────────────────────────────────────
\section{Limitaciones}

\subsection{Cobertura Fija del Prompt}

La base de conocimiento es estática y fue construida manualmente.
No cubre cultivos no incluidos (espárragos, cebollas, berries distintos
al arándano), variedades locales específicas, ni datos dinámicos
(precios de mercado, alertas fitosanitarias activas).
Cada actualización requiere modificar el código y reiniciar el servidor.

\subsection{Capacidad del Modelo (3B parámetros)}

\texttt{llama3.2:3b} es un modelo compacto diseñado para ejecutarse en
hardware de consumo. En comparación con modelos de mayor escala, presenta
limitaciones en razonamiento de múltiples pasos y en la síntesis de
información compleja. En consultas que requieren combinar datos de varias
secciones del prompt simultáneamente se observaron respuestas incompletas
en pruebas no reportadas.

\subsection{Sin Retrieval-Augmented Generation (RAG)}

El sistema no consulta documentos externos en tiempo real.
Agregar RAG con LangChain y ChromaDB permitiría indexar los PDFs
oficiales del SAG e INIA y responder con citación de fuente,
mejorando tanto la cobertura como la verificabilidad de las respuestas.

\subsection{Sin Memoria Persistente}

El historial de conversación se pierde al recargar la página o al
abrir una nueva sesión. No existe persistencia entre sesiones.

% ── 7. Trabajo Futuro ───────────────────────────────────────
\section{Trabajo Futuro}

\textbf{RAG con fuentes oficiales.}
Implementar LangChain\,+\,ChromaDB para indexar documentos del SAG (registro
de plaguicidas), INIA (manuales técnicos) e INDAP (guías de riego).
Esto permitiría actualizar el conocimiento sin modificar el código y
citar la fuente de cada respuesta.

\textbf{Modelos más grandes.}
Migrar a \texttt{llama3.1:8b} (requiere $\approx$8~GB VRAM/RAM) para
mejorar la calidad de respuestas en preguntas de alta complejidad,
manteniendo la misma arquitectura de backend.

\textbf{Evaluación automática.}
Construir un benchmark con preguntas y respuestas de referencia (ground truth)
y medir automáticamente la similitud semántica (ROUGE, BERTScore) entre
respuestas generadas y respuestas esperadas, para facilitar comparaciones
entre versiones del sistema.

\textbf{Clasificador de intención.}
Añadir un módulo de clasificación de consultas que detecte preguntas fuera
del dominio y muestre un mensaje de aclaración, mejorando la experiencia
del usuario y evitando respuestas con alucinaciones en temas no cubiertos.

\textbf{Multimodalidad.}
Permitir subir fotografías de plantas o insectos para diagnóstico visual
usando un modelo multimodal (LLaVA, Llama 3.2 Vision) integrado en
la misma arquitectura Flask.

% ── 8. Conclusiones ─────────────────────────────────────────
\section{Conclusiones}

Se desarrolló AgroBot, un chatbot de asesoría agrícola chilena basado en
prompt engineering y un LLM local de 3B parámetros.
El sistema demuestra que es posible construir un asistente especializado
y funcional con recursos computacionales modestos (8\,GB RAM, sin GPU dedicada),
sin dependencia de APIs externas y con operación offline.

Las 10 preguntas de validación fueron respondidas correctamente
en las áreas de fenología, plagas, riego, normativa laboral, fertilización,
heladas, certificación y cosecha, evidenciando que el prompt engineering
con conocimiento curado es efectivo para dominios técnicos acotados.

La principal limitación es la cobertura fija del sistema, que se aborda
en el trabajo futuro mediante la integración de RAG con documentos oficiales.
En su estado actual, AgroBot representa una prueba de concepto sólida para
el uso de LLMs locales en agricultura, con potencial de escalamiento hacia
un sistema de apoyo a la decisión para pequeños productores chilenos.

% ── Referencias ─────────────────────────────────────────────
\begin{thebibliography}{9}

\bibitem{odepa2023}
ODEPA,
``Estadísticas agrícolas -- Exportaciones frutícolas 2023,''
Oficina de Estudios y Políticas Agrarias, Chile, 2023.
[En línea]. Disponible: \url{https://www.odepa.gob.cl}

\bibitem{singhal2023large}
K.~Singhal \textit{et al.},
``Large language models encode clinical knowledge,''
\textit{Nature}, vol.~620, pp.~172--180, 2023.

\bibitem{huang2023lawyer}
J.~Huang, K.~C.-C. Chang, J.~Gao, A.~Awadallah, and W.~Chen,
``Is ChatGPT a good NLG evaluator? A preliminary study,''
\textit{arXiv preprint arXiv:2303.04048}, 2023.

\bibitem{touvron2023llama}
H.~Touvron \textit{et al.},
``Llama 2: Open foundation and fine-tuned chat models,''
\textit{arXiv preprint arXiv:2307.09288}, 2023.

\bibitem{white2023prompt}
J.~White \textit{et al.},
``A prompt pattern catalog to enhance prompt engineering with ChatGPT,''
\textit{arXiv preprint arXiv:2302.11382}, 2023.

\bibitem{lewis2020rag}
P.~Lewis \textit{et al.},
``Retrieval-augmented generation for knowledge-intensive NLP tasks,''
in \textit{Proc. NeurIPS}, 2020, pp.~9459--9474.

\end{thebibliography}

\end{document}
